% Methods Section for Context Transformation Paper

\section{Methods}

\subsection{Transformation Analysis Framework}

We developed a comprehensive transformation analysis framework to systematically study context-induced changes in token representations. The framework implements 17 complementary analyses, each targeting different aspects of the transformation phenomenon. All code is available at [repository URL].

\subsubsection{Core Infrastructure}

The framework follows a unified pipeline: data loading → validation → analysis → standardized output. Each analysis extends a \texttt{BaseTransformationAnalysis} class ensuring consistent:
\begin{itemize}
\item Data loading with LRU caching for efficiency
\item Validation of inputs and outputs
\item JSON and pickle serialization
\item Comprehensive logging and error handling
\end{itemize}

\subsubsection{Data Preparation}

We analyzed 1,000 tokens from GPT-2's vocabulary, selected to span different frequencies and linguistic types. For each token, we extracted activation trajectories across all 12 layers under 10 different context conditions:
\begin{itemize}
\item Baseline (no context)
\item Determiners: "the", "a"
\item Copula: "is"
\item Modal: "will"
\item Possessive: "my"
\item Conjunction: "and"
\item Negation: "not"
\item Intensifier: "very"
\item Sentence start position
\end{itemize}

Trajectories were obtained by clustering activations at each layer using k-means (k=10), validated through stability analysis across multiple random seeds.

\subsection{Analysis Components}

\subsubsection{Transition Matrix Analysis}
We constructed transition matrices $T_{c,l}$ for each context $c$ and layer $l$, where $T_{c,l}[i,j]$ represents the probability of transitioning from cluster $i$ to cluster $j$. For each matrix, we calculated:
\begin{itemize}
\item Entropy: $H = -\sum_{i,j} T[i,j] \log_2 T[i,j]$
\item Sparsity: $S = 1 - \frac{|\{T[i,j] : T[i,j] > 0\}|}{k^2}$
\item Mutual information between source and destination clusters
\item Diagonal dominance (self-transition probability)
\end{itemize}

\subsubsection{Statistical Validation}
\textbf{Random Baselines:} We compared observed transitions against three random baseline types:
\begin{enumerate}
\item Shuffle: Random permutation of cluster assignments
\item Permute: Random transitions preserving marginal distributions
\item Uniform: Uniformly random transitions
\end{enumerate}

\textbf{Permutation Testing:} For each context and layer, we performed 1,000 permutations to test whether observed patterns differed from chance. P-values were corrected using False Discovery Rate (FDR) with $q < 0.05$.

\textbf{Effect Sizes:} We calculated standardized effect sizes:
\begin{itemize}
\item Cohen's $d$ for continuous metrics
\item Hedge's $g$ for small sample corrections
\item Cliff's $\delta$ for non-parametric comparisons
\end{itemize}

\subsubsection{Information Theory Metrics}
We quantified information content using:
\begin{itemize}
\item Mutual information: $I(Context; Transformation)$
\item KL divergence: $D_{KL}(P_{context} || P_{baseline})$
\item Jensen-Shannon divergence for symmetric comparisons
\item Conditional entropy: $H(Cluster|Context)$
\end{itemize}

\subsubsection{Geometric Analysis}
Procrustes analysis decomposed transformations into geometric components:
\begin{equation}
\min_{R,s,t} ||sXR + t - Y||_F^2
\end{equation}
where $R$ is rotation, $s$ is scaling, and $t$ is translation.

\subsubsection{Predictive Modeling}
We trained machine learning models to predict transformation patterns from token features:
\begin{itemize}
\item Features: token frequency, type, length, linguistic properties
\item Models: Random Forest, Gradient Boosting, SVM
\item Evaluation: 5-fold cross-validation with stratification
\end{itemize}

\subsection{Bootstrap Confidence Intervals}

All metrics include 95\% confidence intervals from 1,000 bootstrap samples. The bootstrap procedure:
\begin{enumerate}
\item Resample tokens with replacement
\item Recalculate metric on bootstrap sample
\item Repeat 1,000 times
\item Use percentile method for CI bounds
\end{enumerate}

\subsection{Computational Details}

Analyses were performed on a system with:
\begin{itemize}
\item GPT-2 (124M parameters) from HuggingFace Transformers
\item PyTorch 2.0 for activation extraction
\item Scikit-learn for clustering and ML models
\item SciPy for statistical tests
\item NumPy/Pandas for data manipulation
\end{itemize}

The complete pipeline processes 1,000 tokens × 10 contexts × 12 layers = 120,000 trajectory points in approximately 2 hours on a single GPU.

\subsection{Visualization}

Publication figures were generated at 300 DPI using Matplotlib with Nature/Science styling guidelines. Interactive visualizations use Plotly for exploration. All figures include:
\begin{itemize}
\item Error bars from bootstrap confidence intervals
\item Consistent color scheme for contexts
\item Clear axis labels and legends
\item Statistical significance indicators where applicable
\end{itemize}